\section{Princípios SOLID}

Os princípios SOLID são um conjunto de diretrizes para o design de software orientado a objetos que visam criar sistemas mais flexíveis, manuteníveis e extensíveis. Estes princípios são:

\begin{itemize}
    \item \textbf{Single Responsibility Principle (SRP)}: Uma classe deve ter apenas uma razão para mudar, ou seja, deve ter apenas uma responsabilidade.
    
    \item \textbf{Open/Closed Principle (OCP)}: Entidades de software devem estar abertas para extensão, mas fechadas para modificação.
    
    \item \textbf{Liskov Substitution Principle (LSP)}: Objetos de uma superclasse devem ser substituíveis por objetos de suas subclasses sem alterar a funcionalidade do programa.
    
    \item \textbf{Interface Segregation Principle (ISP)}: Clientes não devem ser forçados a depender de interfaces que não utilizam.
    
    \item \textbf{Dependency Inversion Principle (DIP)}: Módulos de alto nível não devem depender de módulos de baixo nível. Ambos devem depender de abstrações.
\end{itemize}

Estes princípios são especialmente relevantes para o Modelo de Atores, pois cada ator pode ser visto como uma unidade que segue o princípio de responsabilidade única, enquanto a comunicação baseada em mensagens facilita a inversão de dependências e a segregação de interfaces.
